\section{AutoGen：多 Agent 对话编程框架}

\subsection{核心设计理念}

AutoGen 的设计基于一个核心洞察：\textbf{对话（conversation）是最基础的 Agent 交互原语}。

\begin{knowledgebox}{为什么选择对话作为基础原语}
对话是人类推进学习和证明定理的基本方式。从 ChatGPT 的成功可以看出，多轮对话天然适合迭代改进和协作求解。以对话为中心的设计可以自然地派生出所有其他 Agent 交互模式。
\end{knowledgebox}

构建 Agent 应用只需两步：
\begin{enumerate}
    \item \textbf{定义 Agent}：每个 Agent 可以使用 LLM、工具或人类输入作为后端
    \item \textbf{让 Agent 对话}：通过对话模式（两人对话、嵌套对话、群聊等）编排协作
\end{enumerate}

\subsection{Conversable Agent 抽象}

AutoGen 的统一抽象——\textbf{Conversable Agent}——将不同类型的实体（LLM、工具、人类）统一为可对话的 Agent：
\begin{itemize}
    \item 可以选择 LLM 作为后端（如 GPT-4、Claude）
    \item 可以使用工具（Python 代码执行、API 调用等）
    \item 可以接入人类输入
    \item 可以混合多种后端
    \item 支持嵌套——Agent 内部可以包含其他 Agent 的对话
\end{itemize}

\subsection{对话模式}

\begin{itemize}
    \item \textbf{两人对话}：最简单的模式，已支持 Reflection（写作 Agent + 评审 Agent 迭代改进）
    \item \textbf{Sequential Chat}：多个两人对话串联，如嵌套的多角度评审（SEO 评审 → 法律评审 → 伦理评审 → 元评审）
    \item \textbf{Nested Chat}：Agent 内部发起子对话，如 Commander Agent 先与 Writer 对话获取代码，再与 Safeguard 对话验证安全性
    \item \textbf{Group Chat}：多个不同角色的 Agent 在群组中动态选择发言者，支持约束条件和状态转移逻辑
\end{itemize}

\subsection{Multi-Agent 的优势}

\begin{importantbox}{为什么多 Agent 比单 Agent 好？}
实验表明，将任务分解给专门的 Agent（如 Writer + Safeguard）比让单个 Agent 同时处理所有指令效果更好：GPT-4 多 Agent 方案比单 Agent 方案在安全检查上高出 20\%，GPT-3.5 差距更大。任务越复杂、模型越弱，多 Agent 的优势越明显。
\end{importantbox}

从编程角度看，多 Agent 架构还具有：
\begin{itemize}
    \item 更好的模块化和可维护性
    \item 支持自然的人类参与（人类接管任一 Agent 角色）
    \item 易于扩展（替换单个 Agent 不影响整体）
\end{itemize}

\subsection{AutoBuild：自动构建 Agent 团队}

AutoBuild 解决``为特定任务设计什么 Agent 团队''的问题：
\begin{enumerate}
    \item 用户描述高层需求
    \item 系统自动建议不同角色的 Agent
    \item Agent 组成 Group Chat 协作完成任务
    \item Adaptive Build 进一步支持动态分步构建 Agent 团队
\end{enumerate}

\subsection{本章小结}

AutoGen 以对话为核心原语，提供了灵活且强大的多 Agent 编程框架。其设计哲学是``定义 Agent + 让它们对话''，通过嵌套和组合，可以递归地构建任意复杂的 Agent 系统。

